% Avsnitt 8.1-8.6, ur lectures/L08/appendix/a_state_machines.md (A.1-A.6).

\section{Vad en ändlig tillståndsmaskin är}
\label{c8:sec:what}

En ändlig tillståndsmaskin (FSM) är en krets vars beteende inte bara beror på dess aktuella ingångar
utan på en ändlig mängd \emph{historik}, sammanfattad i ett enda värde som kallas dess
\textbf{tillstånd}. I varje ögonblick befinner sig maskinen i exakt ett tillstånd; vid varje
klockflank kan den gå vidare till ett nytt, bestämt av dess aktuella tillstånd och ingångar; och
dess utgång beräknas ur det tillståndet.

Inget av det här är ny hårdvara. En tillståndsmaskin kombinerar två saker du redan kan: ett
\textbf{register} (\chapref{c3:ch}) som håller det aktuella tillståndet, en vippa per tillståndsbit,
och \textbf{kombinatorik} som beräknar nästa tillstånd och utgången ur det aktuella tillståndet och
ingångarna.

Det nya är disciplinen. I stället för ad hoc-logik som matar en ad hoc-uppsättning vippor arbetar du
i en fast ordning: rita först ett \textbf{tillståndsdiagram}, några namngivna tillstånd med märkta
övergångar mellan sig, och härled tillståndstabellen först när det är rätt, sedan logiken för nästa
tillstånd, sedan grindnätet eller VHDL-koden.

Varje synkron krets hittills är tekniskt sett en trivial tillståndsmaskin: en räknares tillstånd är
dess räkning, ett skiftregisters är dess bitmönster. Det som skiljer en tillståndsmaskin är att
tillståndet varken är en räkning eller ett skiftat mönster utan en godtycklig etikett
(\code{STATE_OFF}, \code{STATE_BLINK}, \dots) vars innebörd du väljer och vars nästa värde får bero
på ingången hur du vill.

Det finns två varianter, och skillnaden avgör vad en utgång får titta på:
\begin{itemize}
  \item \textbf{Mooremaskiner}: utgången beror \emph{bara} på det aktuella tillståndet. Det är den
    sorten som konstrueras för hand här.
  \item \textbf{Mealymaskiner}: utgången beror på det aktuella tillståndet \emph{och} den aktuella
    ingången. \Secref{c8:sec:mealy} går kort igenom dem, bara i VHDL, så att du känner igen en när
    du möter den och kan skillnaden på en klockcykel.
\end{itemize}

% ------------------------------------------------------------------------------------------
\section{Att konstruera en Mooremaskin för hand}
\label{c8:sec:byhand}

Samma process som för varje grindnät i \chapref{c1:ch} och \ref{c2:ch}, tillämpad på en krets med
minne. Det genomarbetade exemplet är en liten, sluten Mooremaskin med fyra tillstånd: en cirkulär
räknare genom en Graykodssekvens, som stegas ett steg per knapptryck.

\subsection*{Specifikation}
\begin{itemize}
  \item Fyra tillstånd, vart och ett representerat av två vippor $\{Q_1, Q_2\}$:

  \begin{narrowtable}{lc}
    \thead{Tillstånd} & \thead{$\{Q_1, Q_2\}$} \\
    \midrule
    \code{STATE_0} & \code{00} \\
    \code{STATE_1} & \code{01} \\
    \code{STATE_2} & \code{11} \\
    \code{STATE_3} & \code{10} \\
  \end{narrowtable}

  \item En systemklocka \code{clock}, och en aktiv låg reset \code{reset_n} som återför maskinen
    till \code{STATE_0}.
  \item En aktiv låg tryckknapp \code{button_n}, ur vilken en synkroniserad, flankdetekterad puls
    \code{X} härleds på exakt det sätt \chapref{c4:ch}:s dubbelvippsynkroniserare gör. \code{X} är
    hög i exakt en klockcykel vid en fallande flank på \code{button_n}.
  \item En utgång \code{Y}, kopplad till en lysdiod, hög bara i \code{STATE_3}.
  \item Maskinen är \textbf{sluten}: tillståndet efter \code{STATE_3} är \code{STATE_0}.
  \item Det är en \textbf{Moore}-maskin: \code{Y} beror bara på $\{Q_1, Q_2\}$, aldrig på \code{X}.
\end{itemize}

Kodningen är en Graykod ($00 \to 01 \to 11 \to 10 \to 00$), där tillstånd som följer på varandra
skiljer sig i exakt en bit. Det är ett medvetet val, inte en slump:
\begin{itemize}
  \item Med en vanlig binärräkning ($00 \to 01 \to 10 \to 11$) vänder en övergång som
    $01 \to 10$ två bitar på en gång, och de två vipporna kommer inte att slå om vid \emph{exakt}
    samma ögonblick.
  \item Den skevheten är harmlös för maskinen själv. Båda vipporna uppdateras vid samma klockflank,
    och ingenting samplar tillståndsregistret däremellan, så maskinen agerar aldrig på
    mellanvärdet. Det är samma argument som \secref{c5:sec:fmax}:s kritiska väg: så länge allt
    hinner stabilisera sig före nästa flank spelar det ingen roll vad som händer under tiden.
  \item Vad skevheten \emph{kan} störa är allt som läser tillståndet kombinatoriskt. En utgång som
    avkodas direkt ur tillståndsbitarna, som \code{Y} nedan, kan glitcha kort medan de är skeva.
    Detsamma gäller ett värde som samplas av en \emph{annan} klockdomän, vilket är fallet
    \secref{c4:sec:sync} varnade för.
  \item En Graykod tar bort båda, eftersom bara en bit någonsin ändras per övergång. Därför spelar
    den störst roll för en räknare som korsar klockdomäner, och minst för en maskin vars utgång är
    registrerad.
\end{itemize}

Kodning är i allmänhet ett designbeslut, och det här är det enda stället i boken där du fattar det
för hand. Binärt använder färrest vippor; Gray minimerar antalet bitar som ändras per övergång;
\textbf{one-hot} lägger en vippa per tillstånd så att logiken för nästa tillstånd blir en enda bred
OR-grind per tillstånd, vilket ofta är snabbast på en FPGA där vippor finns i överflöd. Så snart du
skriver maskinen i VHDL som en uppräknad typ (\secref{c8:sec:vhdl}) lämnar du beslutet till
syntesverktyget, som väljer en kodning åt dig och omkodar fritt om du inte säger något annat.

\subsection*{Tillståndsdiagram}

\lecfig[0.66]{lectures/L08/appendix/images/fsm_state_diagram.png}{Tillståndsdiagram för
Mooremaskinen med fyra tillstånd och Graykod.}{c8:fig:statediagram}

Varje tillståndsbubbla visar det aktuella tillståndet och utgången som \code{Q1Q2/Y}, och varje pil
är märkt med den ingång \code{X} som utlöser den. Varje tillstånd går tillbaka till sig självt så
länge \code{X=0}, och en synkroniserad puls från knappen (\code{X=1}) stegar vidare till nästa
tillstånd och slår runt från \code{STATE_3} tillbaka till \code{STATE_0}.

\subsection*{Tillståndstabell}
Varje kombination av aktuellt tillstånd och ingång, med nästa tillstånd $\{Q_1^{+}, Q_2^{+}\}$ och
utgången \code{Y}:

\begin{narrowtable}{cccccc}
  \thead{$Q_1$} & \thead{$Q_2$} & \thead{$X$} & \thead{$Q_1^{+}$} & \thead{$Q_2^{+}$} &
    \thead{$Y$} \\
  \midrule
  0 & 0 & 0 & 0 & 0 & 0 \\
  0 & 0 & 1 & 0 & 1 & 0 \\
  0 & 1 & 0 & 0 & 1 & 0 \\
  0 & 1 & 1 & 1 & 1 & 0 \\
  1 & 0 & 0 & 1 & 0 & 1 \\
  1 & 0 & 1 & 0 & 0 & 1 \\
  1 & 1 & 0 & 1 & 1 & 0 \\
  1 & 1 & 1 & 1 & 0 & 0 \\
\end{narrowtable}

\subsection*{Karnaughhärledd logik}
Om $Q_1^{+}$ och $Q_2^{+}$ var för sig behandlas som en boolesk funktion av $\{Q_1, Q_2, X\}$ och
körs genom ett Karnaughdiagram, samma teknik som i \chapref{c2:ch}:

\[ Q_1^{+} = Q_1 X' + Q_2 X \qquad Q_2^{+} = Q_2 X' + Q_1' X \]

Lägg märke till formen på båda:
\begin{itemize}
  \item var och en är ”håll kvar när \code{X} är \code{0}, ta något från den andra biten när
    \code{X} är \code{1}”.
  \item hållhalvorna, $Q_1 X'$ och $Q_2 X'$, är verkligen spegelbilder av varandra.
  \item stegningshalvorna är $Q_2 X$ och $Q_1' X$, och komplementet på den andra är inte ett
    skrivfel: det är det som gör sekvensen till en Graykod i stället för en enkel rotation.
\end{itemize}

Utgången \code{Y} förenklas ytterligare, direkt ur tillståndstabellen, utan något beroende av
\code{X} alls, vilket bekräftar att det här är en giltig Mooremaskin:

\[ Y = Q_1 Q_2' \]

\subsection*{Att realisera den som ett grindnät}
$\{Q_1, Q_2\}$ är var sin D-vippa, precis som \chapref{c3:ch}:s register: $Q_1^{+}$ och $Q_2^{+}$
matar D-ingångarna och $Q_1$, $Q_2$ läses tillbaka från Q-utgångarna. \code{Y} är en enda AND-grind
med två ingångar där den ena är inverterad, och den läser det aktuella tillståndet. Det är hela
maskinen: två vippor och en handfull grindar som beräknar $Q_1^{+}$, $Q_2^{+}$ och \code{Y} ur
$\{Q_1, Q_2, X\}$.

% ------------------------------------------------------------------------------------------
\section{Från grindnät till CircuitVerse}
\label{c8:sec:cv}

Innan något av det här skrivs i VHDL, realisera det för hand i CircuitVerse, precis som varje
grindnät i \chapref{c1:ch}: dubbelvippsynkronisera \code{button_n} och \code{reset_n},
flankdetektera den synkroniserade knappen för att producera \code{X}, koppla $Q_1^{+}$ och
$Q_2^{+}$ till ett par D-vippor, och koppla \code{Y} rakt in i lysdioden.

\lecfig[1.0]{lectures/L08/appendix/images/fsm_gate_network.png}{Grindnät som realiserar
Mooremaskinen med fyra tillstånd, inklusive dubbelvippsynkronisering och flankdetektering för
\code{button_n}.}{c8:fig:gatenet}

Läst från vänster till höger:
\begin{itemize}
  \item \code{button_n} går genom tre synkroniseringsvippor
    (\code{button_s1_n}, \code{button_s2_n}, \code{button_s3_n}).
  \item En AND-grind härleder \code{button_edge_s2} (den här kretsens \code{X}) ur andra och tredje
    steget.
  \item De fyra lodräta ledningarna märkta \code{Q1Q2XX'} bär tillståndsbitarna och flankpulsen
    vidare till grindarna till höger, som implementerar \secref{c8:sec:byhand}:s ekvationer för
    $Q_1^{+}$ och $Q_2^{+}$.
  \item De två märkta vipporna lagrar $Q_1$ och $Q_2$ och matar de avslutande AND- och NOT-grindar
    som producerar \code{Y}, kopplad till \code{led}.
  \item \code{reset_n} synkroniseras av sitt eget vippapar nere till vänster, precis som varje annan
    synkroniserad reset i den här boken.
\end{itemize}

Öppna den körbara kretsen själv: \file{lectures/L08/appendix/images/fsm_gate_network.cv}. Stega den
för hand, ett knapptryck i taget, och bekräfta att den besöker
\code{STATE_0} $\to$ \code{STATE_1} $\to$ \code{STATE_2} $\to$ \code{STATE_3} $\to$
\code{STATE_0} i
ordning, med lysdioden tänd bara i \code{STATE_3}.

% ------------------------------------------------------------------------------------------
\section{Samma maskin i VHDL}
\label{c8:sec:vhdl}

Det var sista gången i den här boken du konstruerar en maskin för hand. Allt härifrån skriver den
direkt i VHDL, vilket är värt att ha förtjänat snarare än att ha börjat med: syntesverktyget gör
exakt det du just gjorde på papper, och när en tillståndsmaskin beter sig fel är det diagrammet du
felsöker, inte VHDL-koden.

Översättningen ersätter den explicita tillståndstabellen och Karnaughekvationerna med en
\textbf{uppräknad typ} som namnger varje tillstånd, en \textbf{case-sats} som beskriver varje
tillstånds övergångar, och ett syntesverktyg som lämnas att lista ut vippkodningen och logiken på
grindnivå.

\subsection*{Den uppräknade tillståndstypen}
Anta en maskin med tre tillstånd som styr en lysdiod: \code{STATE_OFF}, \code{STATE_BLINK} (blinkar
i en fast takt) och \code{STATE_ON}. VHDL deklarerar exakt de tre värdena som en ny typ:

\begin{vhdlcode}
type state_t is (STATE_OFF, STATE_BLINK, STATE_ON);
signal state: state_t;
...
state <= STATE_ON;
\end{vhdlcode}

\subsection*{Mönstret med case-satsen}
Maskinen är nästan alltid en synkron process som uppdaterar \code{state} vid stigande flank via ett
\code{case}, en gren per tillstånd, med en asynkron reset som återför den till ett känt
starttillstånd. Antag att \code{to_next_state} och \code{to_prev_state} är synkroniserade,
flankdetekterade pulser producerade på exakt samma sätt som \code{X} i \secref{c8:sec:byhand} och
\ref{c8:sec:cv}:

\begin{vhdlcode}
process(clock, reset_s2_n) is
begin
    if (reset_s2_n = '0') then
        -- On reset, go to STATE_OFF.
        state <= STATE_OFF;
    elsif (rising_edge(clock)) then
        case (state) is
            when STATE_OFF =>
                if (to_next_state = '1') then
                    state <= STATE_BLINK;
                elsif (to_prev_state = '1') then
                    state <= STATE_ON;
                end if;
            -- STATE_BLINK and STATE_ON follow the same shape.
            when others =>
                state <= STATE_OFF;
        end case;
    end if;
end process;
\end{vhdlcode}

\code{case (state) is ... end case;} måste täcka varje värde \code{state_t} kan anta, och
\code{when others} håller det uttömmande om du lägger till ett tillstånd senare och glömmer en gren.

Vad \code{when others} \emph{inte} är, är ett skyddsnät mot ett korrumperat tillståndsregister.
Varje värde \code{state_t} kan anta är redan namngivet ovanför den, så grenen är onåbar i källkoden
och syntesen optimerar normalt bort den i stället för att bygga återhämtningslogik. När Quartus väl
har omkodat maskinen är de flesta bitmönster otillåtna, och att ta sig ur ett av dem kräver Quartus
inställning Safe State Machine snarare än något du kan skriva i VHDL.

\subsection*{Genomarbetat exempel: \code{fsm_led}}
\file{fsm_led.vhd} implementerar exakt den maskinen med tre tillstånd, utökad till två knappar:
\code{button_n(0)} stegar framåt, \code{button_n(1)} går bakåt, tillstånden är slutna i en loop, och
\code{STATE_BLINK} växlar lysdioden var 100:e ms.

Tre delblock återanvänds, alla moduler du skrivit själv:
\begin{itemize}
  \item \code{reset_sync}: \chapref{c4:ch}:s resetsynkroniserare, oförändrad
    (Övning~\ref{c4:ex:resetsync}).
  \item \code{button_sync}: \chapref{c4:ch}:s knappsynkroniserare, också oförändrad, men
    instansierad med \code{generic map(2)} så att den hanterar båda knapparna på en gång och
    producerar en pulsvektor på två bitar. Det är utdelningen från den generic:en: filen är
    identisk med den \chapref{c4:ch} och \ref{c7:ch} instansierar med \code{1}
    (Övning~\ref{c4:ex:buttonsync}).
  \item \code{timer}: \chapref{c7:ch}:s timer, oförändrad (Övning~\ref{c7:ex:timer}).
    \code{fsm_led} deklarerar en egen generic \code{TIMER_TICK_COUNT} och skickar den rakt vidare,
    med förvalet \code{5_000_000} för 100~ms vid \code{50 MHz}, och aktiverar timern bara i
    \code{STATE_BLINK}.
\end{itemize}

Ingen av de tre är utdelad i \file{fsm_led/}: kopiera in dina egna innan du bygger exemplet. Vid det
här laget i boken är det fyra kapitels egna moduler komponerade till en enda design, vilket är det
som är värt att lägga märke till.

\vhdlfile{lectures/L08/fsm_led/fsm_led.vhd}

Två synkrona processer implementerar maskinen:
\begin{itemize}
  \item \code{STATE_PROCESS} är övergångslogiken med \code{case} ovan, driven av
    \code{to_next_state} och \code{to_prev_state}, som härleds kombinatoriskt ur
    \code{button_edge_s2}. Var och en kräver att den \emph{andra} knappens puls är \code{'0'} i
    samma cykel, så att trycka ned båda så nära i tid att deras pulser hamnar på samma klockcykel
    flyttar maskinen ingenstans i stället för att godtyckligt utse en vinnare. En tvetydig ingång
    förtjänar ett definierat svar, och ”gör ingenting” är det som inte kan överraska dig.
  \item \code{LED_PROCESS} driver \code{led_s} enbart ur det aktuella tillståndet: släckt i
    \code{STATE_OFF}, tänd i \code{STATE_ON}, växlande vid varje timeout från timern i
    \code{STATE_BLINK}. Eftersom den aldrig läser \code{button_n} eller \code{button_edge_s2}
    direkt, bara \code{state}, är det här en äkta Mooremaskin: lysdioden reagerar på ett tryck bara
    indirekt, genom det tillstånd trycket orsakade.
\end{itemize}

% ------------------------------------------------------------------------------------------
\section{Mealymaskiner: skillnaden på en klockcykel}
\label{c8:sec:mealy}

Allt hittills har varit en \textbf{Moore}-maskin, med flit: Moore är det säkrare standardvalet och
den form du konstruerar för hand. En \textbf{Mealy}-maskin luckrar upp regeln som definierar en
Mooremaskin: utgången får bero på det aktuella tillståndet \emph{och} den aktuella ingången,
beräknad kombinatoriskt. Den lagras inte i ett register, och den ändras inom samma klockcykel som
ingången gör.

\subsection*{Beteendet: detektera två \code{1}-bitar i följd}
En seriell ingång \code{din} bär en bit per klockcykel, och \code{y} ska indikera att de två senaste
bitarna båda var \code{1}, där överlappande träffar räknas, så \code{111} rapporterar en träff två
gånger.

Som Mealymaskin behöver det här bara två tillstånd, eftersom utgången kan reagera på ingången
\emph{innan} den syns i nästa tillstånd: \code{STATE_IDLE} (föregående bit var \code{0}, eller det
här är den första) och \code{STATE_ONE} (föregående bit var \code{1}).

\begin{narrowtable}{lccc}
  \thead{Tillstånd} & \thead{\code{din}} & \thead{Nästa tillstånd} & \thead{\code{y}} \\
  \midrule
  \code{STATE_IDLE} & 0 & \code{STATE_IDLE} & 0 \\
  \code{STATE_IDLE} & 1 & \code{STATE_ONE}  & 0 \\
  \code{STATE_ONE}  & 0 & \code{STATE_IDLE} & 0 \\
  \code{STATE_ONE}  & 1 & \code{STATE_ONE}  & \textbf{1} \\
\end{narrowtable}

Tillståndsprocessen är \secref{c8:sec:vhdl}:s \code{case}-mönster; bara utgången skiljer sig:

\vhdlfile{lectures/L08/seq_detect_mealy/seq_detect_mealy.vhd}

Tilldelningen av utgången är en enda konkurrent sats, helt utanför varje klockad process:

\begin{vhdlcode}
y <= '1' when (state = STATE_ONE and din = '1') else '0';
\end{vhdlcode}

Det är en \textbf{villkorlig signaltilldelning}, och det är den sista biten VHDL-syntax den här
boken introducerar. Den är en konkurrent sats, så den hör hemma i arkitekturkroppen snarare än inuti
en process, och den beskriver en multiplexer: välj det första värde vars villkor gäller, annars
värdet efter \code{else}. Det är samma hårdvara du skulle få ur ett \code{if}/\code{else} inuti en
kombinatorisk process, skriven på en rad i stället för fem, och den är värd att ha eftersom en
tilldelning av utgången på en rad gör det uppenbart vid en blick vad utgången beror på.

\begin{vhdlcode}
-- These two describe exactly the same multiplexer.
x <= a when sel = '1' else b;

process(sel, a, b) is
begin
    if (sel = '1') then
        x <= a;
    else
        x <= b;
    end if;
end process;
\end{vhdlcode}

Ge den alltid ett \code{else}. Utelämnar du ett har du beskrivit en signal som behåller sitt gamla
värde när inget villkor gäller, vilket är ett lås, precis som i \secref{c5:sec:latch}.
\code{fsm_led} använder den här formen tre gånger, och Övning~\ref{c8:ex:mealy101} ber dig använda
den.

\textbf{Den enda raden är hela skillnaden.} Den läser \code{din}, alltså är den Mealy. Flytta samma
villkor in i den klockade processen, eller ta bort \code{din} ur det, så har du en Mooremaskin i
stället.

\subsection*{Vad den kostar och vad den köper}
Mooreversionen av den här detektorn behöver ett \textbf{tredje} tillstånd, \code{M_TWO}, som finns
enbart för att i en cykel minnas att träffen inträffade, så att utgången har något att läsa som inte
beror på ingången. Mata båda maskinerna med \code{1, 1, 0, 1, 1, 1}:

\begin{narrowtable}{cclcll}
  \thead{Cykel} & \thead{\code{din}} & \thead{Mealytillstånd (före)} & \thead{Mealy \code{y}} &
    \thead{Mooretillstånd (före)} & \thead{Moore \code{y}} \\
  \midrule
  1 & 1 & \code{STATE_IDLE} & 0 & \code{M_IDLE} & 0 \\
  2 & 1 & \code{STATE_ONE}  & \textbf{1} & \code{M_ONE}  & 0 \\
  3 & 0 & \code{STATE_ONE}  & 0 & \code{M_TWO}  & \textbf{1} \\
  4 & 1 & \code{STATE_IDLE} & 0 & \code{M_IDLE} & 0 \\
  5 & 1 & \code{STATE_ONE}  & \textbf{1} & \code{M_ONE}  & 0 \\
  6 & 1 & \code{STATE_ONE}  & \textbf{1} & \code{M_TWO}  & \textbf{1} \\
\end{narrowtable}

Moores \code{y}-kolumn är exakt Mealys \code{y}-kolumn förskjuten en rad nedåt. Det är hela
avvägningen: Mealy rapporterar träffen \emph{samma} cykel som den andra \code{1}:an kommer där Moore
tar en till, och Mealy behöver 2 tillstånd där Moore behöver 3.

\subsection*{Vilken du ska välja}
Moore, nästan alltid, och det är den resten av boken använder, grupprojektet inräknat:
\begin{itemize}
  \item dess utgång beror inte på någon primär ingång, så den kan inte följa en brusig sådan.
  \item en Mealyutgång kan i princip glitcha kombinatoriskt om dess ingångar inte redan är rena.
  \item observera att ”Moore” i sig inte betyder glitchfri. En utgång som avkodas kombinatoriskt ur
    flera tillståndsbitar, som \secref{c8:sec:byhand}:s $Y = Q_1 Q_2'$, kan fortfarande glitcha
    medan de bitarna är skeva vid en övergång. Det som gör en utgång glitchfri är att
    \emph{registrera} den, så som \code{fsm_led}:s \code{LED_PROCESS} gör, och det är vad du vill ha
    på allt som driver en ledning ut från kretsen. En seriell \code{tx}-ledning som avkodas direkt
    ur tillståndsbitarna kommer att sända ut falska flanker.
  \item en maskin vars utgång beror på färre saker är en maskin med färre sätt att bli fel på.
\end{itemize}

Ta till Mealy bara när du specifikt behöver den lägre latensen eller det mindre antalet tillstånd,
och var medveten om att du då byter en synkron utgång mot en kombinatorisk.

% ------------------------------------------------------------------------------------------
\section{På kortet}
\label{c8:sec:board}

\code{fsm_led} syntetiseras och körs precis som varje VHDL-design gjort sedan \chapref{c1:ch}: skapa
ett Quartus Prime Lite-projekt med DE0-CV:s \code{5CEBA4F23C7N} som målkrets, lägg till
\file{.vhd}-filerna för exemplet och dess delblock, kompilera, tilldela varje port en fysisk pinne i
Pin Planner, och kompilera sedan om och programmera kortet.

För \code{fsm_led}, tilldela \code{clock} pinnen för 50 MHz-oscillatorn, \code{reset_n} och
\code{button_n(1 downto 0)} tre tryckknappar, och \code{led} en lysdiod på kortet. Bekräfta sedan att
upprepade tryck på ”nästa” får lysdioden att cykla släckt $\to$ blinkande $\to$ tänd $\to$ släckt,
att ”föregående” går åt andra hållet, och att reset återför den till släckt.

\code{seq_detect_mealy} demonstreras med flit \textbf{inte} på kortet. Den konsumerar en bit per
klockcykel, och en omkopplare som slås om för hand håller sitt värde i miljontals cykler, så det
enda en lysdiod kan visa är det degenererade fallet: \code{din} hållen hög, \code{y} hög från andra
cykeln och framåt. Skillnaden på en klockcykel som hela \secref{c8:sec:mealy} handlar om sker
alldeles för fort för att se. Kör dess testbänk i stället, vilket är den allmänna lärdomen snarare
än ett undantag:

\begin{shell}
cd lectures/L08/seq_detect_mealy
cp ../../L04/exercises/reset_sync/reset_sync.vhd .     # the one you wrote in L04
ghdl -a --std=93 reset_sync.vhd seq_detect_mealy.vhd seq_detect_mealy_tb.vhd
ghdl -e --std=93 seq_detect_mealy_tb
ghdl -r --std=93 seq_detect_mealy_tb --assert-level=error --stop-time=10ms
\end{shell}

Exakta pinnummer beror på vilka knappar och lysdioder på DE0-CV du väljer, tilldelade på samma sätt
som för varje design sedan \chapref{c1:ch}. Slå upp DE0-CV:s egen dokumentation för dess fysiska
pinnkonfiguration om du behöver bekräfta vilken headerpinne som motsvarar vilken märkt knapp,
omkopplare eller lysdiod.
